\section{The generating function}
\label{dist:sec:pgf}

Everything in this chapter is read off one object. The PGF of the load,
\begin{equation}
G(z,t)=\sum_{n\ge0}p_n(t)z^n,\qquad p_n(t)=\pr{\Xt=n},
\label{dist:eq:gdef}
\end{equation}
packs the entire state distribution into a single function of two variables
and satisfies a first-order partial differential equation that can be solved
exactly. We derive that equation, solve it, and record what the two rupture
conventions do to the answer.

One feature of the catastrophe convention should be named before it is used.
Ruptured realisations are absorbed into the separate state $H$ and no longer
contribute to the count, so $G(1,t)=I(t)\le1$: the PGF is \emph{defective},
and its mass decays as cells burst. Under killing the rupture transition lands
in state $0$ and is counted, so $G(1,t)=1$ for all $t$.

\subsection{From the master equations to a PDE}
\label{dist:sec:masterpde}

The forward Kolmogorov equations for the killing convention are
\begin{equation}
\ddt{p_i} = \lambda (i-1) p_{i-1} + \mu (i+1) p_{i+1}
- (\lambda + \mu + \delta)ip_i \quad (i>0),
\qquad
\ddt{p_0} = \mu p_1 + \delta\sum_{j\ge1} jp_j.
\label{dist:eq:kolmogorov}
\end{equation}
The summation collects every rupture transition --- from any positive state a
collapse to $0$ can occur --- and is itself the mean $\E{\Xt}=J(t)$.
Under catastrophe that term is absent, since rupture leaves the count
altogether; the first equation then holds for all $i\ge0$ with $p_{-1}=0$, and
at $i=0$ reads
\begin{equation}
\ddt{p_0}=\mu p_1,
\label{dist:eq:kolmogorovcat}
\end{equation}
internal extinction by death still feeding state $0$.

Multiply the $i$th equation by $z^i$ and sum over $i$. Shifting indices in the
three sums and using
$\sum_{i\ge1}ip_iz^{i-1}=\pdd{}{z}\sum_{i\ge1}p_iz^{i}=\pdd{G}{z}$ collapses
the infinite system into one equation. Under killing,
\begin{equation}
\pdd{G}{t} = \delta J(t) +
\bigl(\lambda z^2 - (\lambda + \mu + \delta)z + \mu\bigr)\pdd{G}{z},
\label{dist:eq:killpde}
\end{equation}
and dropping the inhomogeneous term --- no rupture flux enters any state of
the count under catastrophe --- leaves
\begin{equation}
\pdd{G}{t} =
\bigl(\lambda z^2 - (\lambda + \mu + \delta)z + \mu\bigr)\pdd{G}{z}.
\label{dist:eq:catpde}
\end{equation}

\subsection{Solving by characteristics}
\label{dist:sec:characteristics}

\Cref{dist:eq:catpde} is first order and linear, so it reduces to a family of
ordinary differential equations along characteristic curves. The generic
construction --- the change of variables $(z,t)\to(s,\xi)$, the admissible
initial curves, and the transport equations
$\dda{G}{\xi}=0$ and $\dda{t}{\xi}=-1$ --- is set out in
\mextref{m:sec:moc}{the method of characteristics in \ChM},
with the generic system at
\mextref{m:eq:genericcharacteristics}{the generic characteristic equations};
we take it from
there and record only what is specific to this process. With the initial curve
$z(s,0)=s$, $t(s,0)=0$, transport gives $t=-\xi$ and $G$ constant on curves of
fixed $s$, so that the initial condition $G(z,0)=z$ --- a single founder,
$p_1(0)=1$ --- fixes $G(s,\xi)=s$.

That leaves one equation, and it is the one that is not generic:
\begin{equation}
\dda{z}{\xi} = \lambda z^2 - (\lambda + \mu + \delta)z + \mu.
\label{dist:eq:zchar}
\end{equation}
This is the Riccati equation \cref{dist:eq:riccati} of the no-rupture
probability, in a different variable and with a different initial condition:
$I(0)=1$ there, $z(s,0)=s$ here. Its solution therefore reproduces the
formula for $I$ with $s$ wherever a $1$ stood,
\begin{equation}
z = \frac{a(s-b) - b(s-a)\ee^{\theta\xi}}{(s-b) - (s-a)\ee^{\theta\xi}},
\label{dist:eq:zsolve}
\end{equation}
now in $\xi$ rather than $t$. Since $G(z,t)=s(z,t)$, one step finishes the
calculation: rearrange for $s$ and put $-t$ for $\xi$. The map is an
involution --- $s$ and $z$ exchange places and nothing else moves --- so
\begin{equation}
s = \frac{a(z-b) - b(z-a)w}{(z-b) - (z-a)w},
\qquad w=\ee^{\theta t},
\label{dist:eq:sinv}
\end{equation}
and the generating function of the single-founder process is
\begin{equation}
G(z,t) = \frac{a(z-b) - b(z-a)w}{(z-b) - (z-a)w}
= \frac{ab(1-w) - z(a-bw)}{(b-aw) - z(1-w)}.
\label{dist:eq:gcat}
\end{equation}
The second form is the one to differentiate: its $z$-dependence is one linear
factor over another, so repeated $\partial_z$ returns a power over a power,
and \cref{dist:sec:geometric} reads the state probabilities off it.

\subsection{Several founders}
\label{dist:sec:pgf-k-founders}

The solution extends to any initial load by the branching property. If the
cell begins with $k$ independent founders, the $k$ lineages evolve
independently until the shared catastrophe clock rings and the total count is
their sum, so the PGF is the $k$th power of the single-founder one:
\begin{equation}
G_k(z,t) = \bigl(G(z,t)\bigr)^k .
\label{dist:eq:gk}
\end{equation}
Under catastrophe this is defective, with $G_k(1,t)=I(t)^k$; under killing it
remains proper. \Cref{dist:eq:gk} is what
\cref{dist:sec:moi,dist:sec:chained} are built on, and it is where the
independence assumption of \cref{dist:sec:assumptions} does its work.

\subsection{The two conventions differ by a function of time}
\label{dist:sec:conventions}

The killing convention needs no second derivation. Its PDE
\cref{dist:eq:killpde} carries the inhomogeneous term $\delta J(t)$, which is
$-\ddt{I}$; along a characteristic that term integrates to $I$ alone, with no
$z$ in it. The conclusion is exact and immediate.

\begin{proposition}[The two rupture conventions]
\label{dist:prop:conventions}
For every $z$ and every $t$,
\begin{equation}
G_{\rm kill}(z,t)=G_{\rm cat}(z,t)+\bigl(1-I(t)\bigr).
\label{dist:eq:conventions}
\end{equation}
\end{proposition}

\begin{proof}
Along the characteristics of \cref{dist:sec:characteristics} the killing
system differs from the catastrophe system in its first equation alone, which
becomes $\dda{G}{\xi}=-\delta J(-\xi)=\dda{}{\xi}I(-\xi)$, since
$\dd t=-\dd\xi$ and $I'=-\delta J$. Hence $G(s,\xi)=-I(-\xi)+\phi(s)$, and
$G(s,0)=s$ with $I(0)=1$ fixes $\phi(s)=1+s$, so $G(s,\xi)=1-I(-\xi)+s$. The
remaining characteristic equation is \cref{dist:eq:zchar} unchanged, giving
\cref{dist:eq:sinv} again; substituting yields
\begin{equation}
G_{\rm kill}(z,t)=1-I(t)+\frac{ab(1-w)-z(a-bw)}{(b-aw)-z(1-w)},
\label{dist:eq:gkill}
\end{equation}
which is \cref{dist:eq:gcat} plus $1-I(t)$.
\end{proof}

Three consequences follow in a line each. The state probabilities agree,
$p_n^{\rm kill}=p_n^{\rm cat}$ for every $n\ge1$. The difference is confined
to the empty state, $p_0^{\rm kill}=p_0^{\rm cat}+(1-I)$. Every moment of
$\Xt$ restricted to $n\ge1$ is therefore convention-independent, so nothing in
the rest of this chapter would change had killing been chosen instead. One
reason to prefer catastrophe nonetheless survives the equivalence: it keeps
$G(1,t)=I(t)$, so that the defect of the generating function is itself the
biological observable, the probability that the cell has not yet ruptured.
That is why the convention was chosen, and why the parallel track need not be
carried further.

With the convention settled, the generating function can be read for what it
contains. Extracting its coefficients is mechanical; what comes out of them is
not.
